% Audited ultraviolet-limit proof; insert before the source section.
% Uses only the host document's packages and theorem environments.
\section{The ultraviolet limit of the flat Gaussian baseline}
\label{limit-section}

The common complex strip gives more than uniform locality: the flat
precision and covariance have limits as the original spacing tends to
zero. We prove a quantitative rate, including the aliases near the edge
of the fine Brillouin zone. This is a limit of the homogeneous free
baseline in the fixed retained coordinates.

For a matrix kernel $F(i,j)$ on a lattice with finite-dimensional fibers,
use the exponentially weighted Schur norm
\begin{equation}\label{limit-schur-definition}
 \norm{F}_\kappa=
 \max\left\{
 \sup_i\sum_j e^{\kappa\operatorname{dist}(i,j)}\norm{F(i,j)},
 \sup_j\sum_i e^{\kappa\operatorname{dist}(i,j)}\norm{F(i,j)}
 \right\}.
\end{equation}
The fiber norm is the fixed operator norm. On the unit lattice the
distance is $|i-j|_1$. For reduced operators we group all coordinates
into the fixed $L$-blocks and use block distance. Grouping a unit-lattice
kernel changes constants and, if needed, reduces its exponent; these
changes depend only on $d,L$ and the fixed fiber dimension. Matrix
products below use one common grouping. The triangle inequality for
distance proves
\begin{equation}\label{limit-schur-algebra}
 \norm{FG}_\kappa\le\norm{F}_\kappa\norm{G}_\kappa.
\end{equation}
The same convention applies in periodic distance on a torus, with the
actual periodic incidence multiplicities.

\begin{theorem}\label{limit-theorem}
There are a periodic holomorphic symbol $q_\infty$ and constants
$\rho_*,C>0$, independent of odd $M$, such that
\begin{equation}\label{limit-symbol-rate}
 \sup_{|\operatorname{Im}z|_\infty\le\rho_*}
       \norm{q_M(z)-q_\infty(z)}\le CM^{-2}.
\end{equation}
At a fixed sufficiently small $\kappa>0$,
\begin{equation}\label{limit-action-rate}
 \norm{q_M-q_\infty}_\kappa\le C_\kappa M^{-2}.
\end{equation}
Set
\[
 A_{\infty,0}=K_L^*(q_\infty+T_L^*T_L)K_L,
 \qquad C_{M,0}=A_{M,0}^{-1},\quad
 C_{\infty,0}=A_{\infty,0}^{-1}.
\]
Then $A_{\infty,0}\ge\gamma_*I$, and at a possibly smaller fixed
positive exponent,
\begin{equation}\label{limit-covariance-rate}
 \norm{A_{M,0}-A_{\infty,0}}_\kappa+
 \norm{C_{M,0}-C_{\infty,0}}_\kappa
 \le C_{d,L,\kappa}M^{-2}.
\end{equation}
The same bounds hold uniformly on compatible periodic tori. Their
kernels are exact periodizations of the infinite kernels. If the
shortest nonzero period tends to infinity, the thermodynamic and
$M\to\infty$ limits commute on every fixed finite set of coordinates.
\end{theorem}

\subsection{Sinc notation and the exact limiting alias sum}

Continue to use $I_M=\{-(M-1)/2,\ldots,(M-1)/2\}$ and
$p_r(z)=z+2\pi r$. To distinguish finite and limiting denominators,
write the earlier $P_r$ as $P_{M,r}$. Define
\begin{equation}\label{limit-sinc}
 s_M(w)=\frac{\sin(w/2)}{M\sin(w/(2M))},\qquad
 s(w)=\frac{\sin(w/2)}{w/2},\qquad s(0)=1,
\end{equation}
with all removable values filled, and
\begin{equation}\label{limit-denominators}
 P_{M,r}=4M^2\sum_\nu\sin^2\frac{p_{r,\nu}}{2M},
 \qquad P_{\infty,r}=\sum_\nu p_{r,\nu}^2.
\end{equation}
The paired averaging numerators are
\begin{align}
 N_{M,r,\mu}
 &=s_M(p_{r,\mu})^2\prod_\nu s_M(p_{r,\nu})^2,
 \label{limit-numerator-finite}\\
 N_{\infty,r,\mu}
 &=s(p_{r,\mu})^2\prod_\nu s(p_{r,\nu})^2.
 \label{limit-numerator-infinite}
\end{align}
Indeed the longitudinal averaging phase
$\exp(i(M-1)p_\mu/(2M))$ tends to $\exp(ip_\mu/2)$, but it
cancels exactly between the positive and negative arguments of each
covariance numerator. There is consequently no phase error of order
$M^{-1}$. These are analytic products, not complex absolute squares.
Physical-volume Parseval normalization gives no additional power of
$M$ in the alias sum.

Define the high-alias remainders
\begin{equation}\label{limit-remainders}
 R_{M,\mu}=
 \sum_{r\in I_M^d\setminus\{0\}}\frac{N_{M,r,\mu}}{P_{M,r}},
 \qquad
 R_{\infty,\mu}=
 \sum_{r\in\mathbb Z^d\setminus\{0\}}
                         \frac{N_{\infty,r,\mu}}{P_{\infty,r}}.
\end{equation}
We first prove normal convergence and a uniform rate for the second
sum. These estimates also justify the explicit continuum formula below.

\subsection{High-alias bounds uniform up to the fine cutoff}

Fix $\pi<K<2\pi$ and work on
\[
 \mathcal B_{K,\rho}=
 \{x+iy:|x_\nu|\le K,\ |y_\nu|\le\rho\},
 \qquad 0<\rho\le1.
\]
Put $a=2\pi-K$ and $c_0=4a^2/\pi^2$. If a retained component
$r_\nu\ge1$, the two possible nearest distances of
$x_\nu+2\pi r_\nu$ to $2\pi M\mathbb Z$ obey
\[
 2\pi r_\nu-K\ge ar_\nu,\qquad
 2\pi(M-r_\nu)-K
 \ge2\pi(r_\nu+1)-K\ge ar_\nu.
\]
The negative case is symmetric. The principal distance is at most
$\pi M$, so the elementary sine lower bound gives
\[
 P_{M,r}(x)\ge c_0|r|_2^2\qquad(r\ne0).
\]
Moreover, direct expansion of the complex sine gives
\[
 \operatorname{Re}P_{M,r}(x+iy)
 \ge P_{M,r}(x)-e^\rho\sum_\nu y_\nu^2.
\]
Choose $\rho$ so that $de^\rho\rho^2\le c_0/2$. Then
\begin{equation}\label{limit-high-denominator}
 |P_{M,r}(z)|\ge\operatorname{Re}P_{M,r}(z)
 \ge\frac{c_0}{2}|r|_2^2\qquad(r\ne0).
\end{equation}
This bound includes $|r_\nu|$ comparable to $M$. The continuum
denominator, for every nonzero integer alias, satisfies
\begin{equation}\label{limit-high-denominator-infinite}
 \operatorname{Re}P_{\infty,r}(z)
 =\sum_\nu(x_\nu+2\pi r_\nu)^2-\sum_\nu y_\nu^2
 \ge\frac{c_0}{2}|r|_2^2.
\end{equation}
For this last inequality use $|x_\nu+2\pi r_\nu|\ge a|r_\nu|$
when $r_\nu\ne0$, and $a^2\ge c_0$.

The integral representation of sinc bounds the zero-index factor by
$e^{\rho/2}$. For $r_\nu\ne0$,
\[
 |\sin((z_\nu+2\pi r_\nu)/2)|\le\cosh(\rho/2),
 \qquad |z_\nu+2\pi r_\nu|\ge a|r_\nu|.
\]
Hence
\begin{equation}\label{limit-sinc-weight}
 |s(p_{r,\nu})|\le\frac{C_s}{1+|r_\nu|}.
\end{equation}
The same estimate holds for retained $s_M$ factors. For a nonzero
index, use the preceding distance estimate to obtain
\[
 |M\sin(p_{r,\nu}(z)/(2M))|
 \ge |M\sin((x_\nu+2\pi r_\nu)/(2M))|
 \ge(a/\pi)|r_\nu|;
\]
for zero index, use the centered finite exponential average.

Set
\[
 \mathcal W_{r,\mu}=(1+|r_\mu|)^{-2}
                         \prod_\nu(1+|r_\nu|)^{-2},
 \qquad
 B_2=\sum_{m\in\mathbb Z}(1+|m|)^{-2}=\pi^2/3-1.
\]
Equations \eqref{limit-high-denominator}--\eqref{limit-sinc-weight}
give the summable majorant
\begin{equation}\label{limit-majorant}
 \left|\frac{N_{M,r,\mu}}{P_{M,r}}\right|
 +\left|\frac{N_{\infty,r,\mu}}{P_{\infty,r}}\right|
 \le C\frac{\mathcal W_{r,\mu}}{|r|_2^2}.
\end{equation}
The finite term here is used only for retained aliases. In particular,
the infinite remainder is holomorphic in the box interior and
continuous uniformly on the closed box. Its tail has the explicit bound
\begin{equation}\label{limit-tail}
 \sum_{|r|_\infty>N}\frac{\mathcal W_{r,\mu}}{|r|_2^2}
 \le\frac{2d B_2^{d-1}}{(N+1)^3}.
\end{equation}
To verify it, discard the additional orientation factor, use
$|r|_2\ge N+1$, and cover the tail by the $d$ coordinate tails;
$\sum_{|m|>N}(1+|m|)^{-2}\le2/(N+1)$.

\subsection{A second-order rate without a small-alias restriction}

Let $f(w)=w/\sin w$, with its removable value at zero. For all
retained aliases and $M\ge3$,
\begin{equation}\label{limit-compact-rectangle}
 \left|\operatorname{Re}\frac{p_{r,\nu}}{2M}\right|
 \le\frac\pi2+\frac{K-\pi}{6}<\frac{2\pi}3,
 \qquad
 \left|\operatorname{Im}\frac{p_{r,\nu}}{2M}\right|\le\frac16.
\end{equation}
Thus $f$ and $(f(w)-1)/w^2$ are bounded on one fixed rectangle,
after filling the removable value at zero. The exact identity
\begin{equation}\label{limit-factor-comparison}
 s_M(p)=s(p)f(p/(2M)),\qquad
 |f(p/(2M))-1|\le C|p|^2/M^2
\end{equation}
holds even at the largest aliases. For $r\ne0$,
$\sum_\nu|p_{r,\nu}|^2\le C|r|_2^2$. Telescoping the fixed
$2d+2$ factors in the numerator therefore gives
\begin{equation}\label{limit-numerator-error}
 |N_{M,r,\mu}-N_{\infty,r,\mu}|
 \le CM^{-2}|r|_2^2\mathcal W_{r,\mu}.
\end{equation}
The removable analytic function
$((\sin w/w)^2-1)/w^2$ is bounded on the same rectangle, yielding
\begin{equation}\label{limit-denominator-error}
 |P_{M,r}-P_{\infty,r}|
 \le CM^{-2}\sum_\nu|p_{r,\nu}|^4
 \le CM^{-2}|r|_2^4.
\end{equation}
Divide \eqref{limit-numerator-error} by the lower denominator bound.
For the reciprocal-denominator difference, divide
\eqref{limit-denominator-error} by the product of the two lower
bounds. The powers of $|r|$ cancel in both operations, giving
\begin{equation}\label{limit-summand-error}
 \left|\frac{N_{M,r,\mu}}{P_{M,r}}
       -\frac{N_{\infty,r,\mu}}{P_{\infty,r}}\right|
 \le CM^{-2}\mathcal W_{r,\mu}.
\end{equation}
Since $\sum_r\mathcal W_{r,\mu}\le B_2^d$, these retained-alias
errors sum to $O(M^{-2})$. The missing aliases are bounded by
\eqref{limit-tail} with $N=(M-1)/2$, giving $O(M^{-3})$. Consequently
\begin{equation}\label{limit-remainder-rate}
 \sup_{z\in\mathcal B_{K,\rho}}
 |R_{M,\mu}(z)-R_{\infty,\mu}(z)|\le CM^{-2}.
\end{equation}
No estimate above assumes $|r|/M$ is small. The case $M=1$ can be
included by enlarging the constant; it has no high aliases.

\subsection{The canceled continuum symbol and its removable zeros}

To avoid confusion with the reduced precision $A_{M,0}$, denote the
principal scalar averaging weight in this subsection by
$\mathfrak a_M$. Define its limit and the other limiting quantities by
\begin{align}
 P(z)&=\sum_\nu z_\nu^2,&
 \mathfrak a(z)&=\prod_\nu s(z_\nu)^2,&
 g_{\infty,\mu}(z)&=\mathfrak a(z)s(z_\mu)^2,
 \label{limit-central}\\
 D_{\infty,\mu}&=g_{\infty,\mu}+P R_{\infty,\mu},&
 h_{\infty,\mu}&=D_{\infty,\mu}^{-1}.&&
 \nonumber
\end{align}
Here $\mathfrak a_M(z)=\prod_\nu s_M(z_\nu)^2$ is the quantity
previously called the central $A(z)$, not a retained-coordinate matrix.
Keep $d_\mu^\pm=e^{\pm iz_\mu}-1$ and $b_\mu=d_\mu^+d_\mu^-$,
and set
\begin{equation}\label{limit-E}
 E_\infty=\mathfrak a^{-1}-\sum_\mu
 \frac{b_\mu R_{\infty,\mu}}
      {g_{\infty,\mu}D_{\infty,\mu}}.
\end{equation}
The explicit limiting precision is
\begin{equation}\label{limit-explicit-symbol}
 (q_\infty)_{\mu\nu}(z)
 =P(z)h_{\infty,\mu}(z)\delta_{\mu\nu}
 -\frac{h_{\infty,\mu}(z)d_\mu^+(z)d_\nu^-(z)
                  h_{\infty,\nu}(z)}{E_\infty(z)}.
\end{equation}

On real nonzero momenta, this is the gauge quotient of the continuum
alias covariance with diagonal entries
\begin{equation}\label{limit-continuum-covariance}
 c_{\infty,\mu}(x)=
 \sum_{r\in\mathbb Z^d}
 \frac{s(x_\mu+2\pi r_\mu)^2
               \prod_\nu s(x_\nu+2\pi r_\nu)^2}
      {\sum_\nu(x_\nu+2\pi r_\nu)^2}.
\end{equation}
Explicitly, if $W_\infty=\diag(c_{\infty,\mu})$, it equals
$W_\infty^{-1}-W_\infty^{-1}d^+
(d^-W_\infty^{-1}d^+)^{-1}d^-W_\infty^{-1}$.
The singular principal Feynman term of
\eqref{limit-continuum-covariance} is not used at $x=0$; the canceled
formula \eqref{limit-explicit-symbol} supplies the continuation.

Choose a closed smaller box inside the established common holomorphic
domain, still with a real margin beyond $[-\pi,\pi]^d$. The finite
denominators $g_{M,\mu},D_{M,\mu},E_M$ have uniform positive modulus
lower bounds there. At the central alias,
\eqref{limit-factor-comparison} and the direct sine expansion give
$O(M^{-2})$ convergence of $P_{M,0}$ to $P$ and
$\mathfrak a_M,g_{M,\mu}$ to their limits. Together with
\eqref{limit-remainder-rate}, this gives $O(M^{-2})$ convergence of
$D_{M,\mu}$ and then $E_M$. Their lower bounds pass to the limits.
The canceled rational expression is uniformly Lipschitz on this
bounded set separated from its denominator zeros. This proves
\eqref{limit-symbol-rate} on the smaller strip.

Every one-variable zero at $z_\nu+2\pi r_\nu=0$ is filled by the
entire sinc function. If $z_\nu=0$ and $r_\nu\ne0$, that sinc
factor is zero and the alias term vanishes; its high denominator is
still nonzero by \eqref{limit-high-denominator-infinite}. At $z=0$,
all high numerators vanish, while
\[
 \mathfrak a(0)=g_{\infty,\mu}(0)
 =D_{\infty,\mu}(0)=E_\infty(0)=1,
 \qquad P(0)=d^\pm(0)=0.
\]
Thus $q_\infty(0)=0$. The exact identity
$\sum_\mu b_\mu/g_{\infty,\mu}=P/\mathfrak a$ holds with
the removable one-variable factors filled. It removes the whole
complex null cone $P=0$ from \eqref{limit-explicit-symbol}.
There is no exceptional rule on coordinate axes.

The separated central and remainder pieces need not be periodic.
The full finite-$M$ quotients are periodic on their common tube,
so their uniform limit is periodic. Alternatively, reindex the
absolutely convergent full sum \eqref{limit-continuum-covariance}
on the real overlaps and apply the identity theorem. This also
identifies the formula across Brillouin-chart seams.

\subsection{Weighted convergence, inversion, and compatible tori}

Choose $0<\kappa<\rho_0<\rho_*$. In the Fourier convention used
above, shift the $\nu$th contour by $i\rho_0\operatorname{sign}(x_\nu)$.
Equation \eqref{limit-symbol-rate} gives
\[
 \norm{(q_M-q_\infty)(x)}
 \le CM^{-2}e^{-\rho_0|x|_1}.
\]
Summation proves \eqref{limit-action-rate}, with the explicit factor
\begin{equation}\label{limit-weighted-rate}
 \norm{q_M-q_\infty}_\kappa
 \le\frac{C}{M^2}
 \left(\frac{1+e^{-(\rho_0-\kappa)}}
            {1-e^{-(\rho_0-\kappa)}}\right)^d.
\end{equation}
Grouping into the fixed $L$-block fibers preserves this rate.
Because $K_L,T_L$ are fixed finite-range maps,
\[
 A_{M,0}-A_{\infty,0}=K_L^*(q_M-q_\infty)K_L,
 \qquad
 \norm{A_{M,0}-A_{\infty,0}}_\kappa
 \le\norm{K_L}_\kappa^2\norm{q_M-q_\infty}_\kappa.
\]
This is an operator-norm limit as well. Hence the real lower gap
$A_{M,0}\ge\gamma_*I$ passes to $A_{\infty,0}$.
The uniform locality and gap argument already proved for the flat
covariance applies to the limit too. Reducing the exponential weight
below the pointwise inverse-kernel decay exponent yields
\[
 \sup_M\norm{C_{M,0}}_\kappa+\norm{C_{\infty,0}}_\kappa<\infty.
\]
The exact resolvent identity and \eqref{limit-schur-algebra} then give
\[
 C_{M,0}-C_{\infty,0}
 =C_{M,0}(A_{\infty,0}-A_{M,0})C_{\infty,0},
\]
\[
 \norm{C_{M,0}-C_{\infty,0}}_\kappa
 \le\norm{C_{M,0}}_\kappa
     \norm{A_{M,0}-A_{\infty,0}}_\kappa
     \norm{C_{\infty,0}}_\kappa
 \le C_{d,L,\kappa}M^{-2}.
\]
This completes the infinite-volume estimate. In the sequence $M=L^k$
its rate is $O(L^{-2k})$.

Let $\Gamma$ be a compatible full-rank period lattice. Unit-lattice
symbols sampled on the torus Fourier grid have periodized kernels, so
\begin{equation}\label{limit-periodization}
 (q_M^\Gamma-q_\infty^\Gamma)(x,y)
 =\sum_{\gamma\in\Gamma}(q_M-q_\infty)(x-y+\gamma).
\end{equation}
The periodic distance is no greater than the distance of each image;
thus its weighted Schur norm is bounded by
\eqref{limit-weighted-rate}, independently of volume.

For retained operators require $\Gamma\subset L\mathbb Z^d$ and
the same periodic pivot convention. In $L$-block coordinates all
operators have fixed finite-dimensional fibers. Periodization
commutes with multiplication of absolutely summable block kernels:
expand both image sums and decompose the intermediate point into a
fundamental-domain point and its period translate. Therefore
periodizing $A_{M,0}C_{M,0}=I=C_{M,0}A_{M,0}$ proves
\[
 C_{M,0}^\Gamma=(A_{M,0}^\Gamma)^{-1},
\]
with $A_{M,0}^\Gamma$ exactly the finite periodic source matrix.
The same reasoning holds for the limit. This verifies the
volume-uniform version of \eqref{limit-covariance-rate}, including
short tori with coincident images and their correct multiplicities.

If the shortest nonzero period tends to infinity, uniform exponential
decay makes the image error vanish on every fixed finite coordinate
set. Combining this error with the volume-independent $M^{-2}$ bound
proves that the two limits commute. Growing volume with a bounded
periodic direction instead produces a cylindrical limit. The proof
uses exact homogeneous periodization; it says nothing about an
arbitrary Dirichlet or retained-history covariance being a compression
or periodization of the homogeneous one.

\begin{remark}\label{limit-scope}
The limit is the order-zero homogeneous Gaussian baseline. Under the
finite-level constrained-minimizer identification, it is expressed in
the actual retained integration coordinates of the source step. It
does not establish an interacting Yang--Mills ultraviolet limit,
nonflat-background convergence, independent-current derivative
bounds, or arbitrary retained-history boundary matching. Those
interfaces are unchanged, and no sharp cutoff domain is moved by
the calculation.
\end{remark}
